\Subhead{3.6.1}{Марковские и диффузионные процессы. Уравнение Фоккера--Планка. Броуновское движение}

Марковское свойство: при известном текущем состоянии условное распределение будущего не зависит от более ранней истории. Это не означает независимости соседних значений. Переходные плотности удовлетворяют уравнению Чепмена–Колмогорова: переход за два интервала складывается интегрированием по промежуточному состоянию.

Диффузионный процесс имеет непрерывные траектории и малые условные приращения со средним $a(x,t)\dd t$ и дисперсией $b^2(x,t)\dd t$. В смысле Ито
\begin{equation}
 \dd X=a(X,t)\dd t+b(X,t)\dd W,\qquad \langle(\dd W)^2\rangle=\dd t.
\end{equation}
Для пробной функции $f$ формула Ито даёт $\dd\langle f\rangle/\dd t=\langle af'+b^2f''/2\rangle$. Интегрируя по частям и перенося производные с $f$ на плотность, получаем
\begin{equation}\label{eq:fokkerplanck}
 \partial_tp=-\partial_x(ap)+\tfrac12\partial_x^2(b^2p).
\end{equation}
При мультипликативном шуме интерпретации Ито и Стратоновича различаются дрейфовой поправкой; уравнение нельзя выбирать без определения модели шума.

Броуновская скорость подчиняется $m\dd v=-bv\dd t+\sqrt{2bk_BT}\dd W$. Стационарная плотность максвелловская, $\langle v^2\rangle=k_BT/m$, корреляция $(k_BT/m)\e^{-b|\tau|/m}$. На временах $t\gg m/b$ координата диффундирует: $\langle[x(t)-x(0)]^2\rangle\simeq2Dt$, $D=k_BT/b$. Это соотношение Эйнштейна; белая случайная сила и непрерывная координата описывают разные процессы.
